% ==============================================================================
% Methods
% ==============================================================================

\section{Methods}
\label{sec:methods}

\subsection{Models}
\label{sec:models}

We study four models spanning two architecture families and two training
regimes (Table~\ref{tab:models}).  The headline comparison is between
Llama-3.1-8B-Instruct \citep{llama31} and DeepSeek-R1-Distill-Llama-8B
\citep{deepseekr1}: identical architectures (32 transformer layers, 32
query heads with 8 KV heads via grouped-query attention, 4096 hidden
dimensions) differing only in that the latter was distilled from a
reasoning-trained teacher.  Any representational differences between these
two models are attributable to reasoning distillation training, not
architectural confounds.  We include Gemma-2-9B-IT \citep{gemma2} and
R1-Distill-Qwen-7B \citep{deepseekr1} as cross-architecture controls.

\begin{table}[t]
\centering
\caption{Models used in this study.  The Llama pair (rows 1 and 3) shares
  identical architecture, enabling controlled comparison.}
\label{tab:models}
\vspace{4pt}
\small
\begin{tabular}{lccccc}
\toprule
\textbf{Model} & \textbf{Type} & \textbf{Layers} & \textbf{Heads} & \textbf{Hidden} & \textbf{SAE source} \\
\midrule
Llama-3.1-8B-Instruct & Standard & 32 & 32 & 4096 & Llama Scope \\
Gemma-2-9B-IT & Standard & 42 & 16 & 3584 & Gemma Scope \\
R1-Distill-Llama-8B & Reasoning & 32 & 32 & 4096 & Llama Scope\textsuperscript{$\dagger$} \\
R1-Distill-Qwen-7B & Reasoning & 28 & 28 & 3584 & Custom \\
\bottomrule
\end{tabular}
\vspace{2pt}
\raggedright
\footnotesize{\textsuperscript{$\dagger$}Llama Scope SAEs were trained on the
base model; we verify reconstruction fidelity on the distilled model before
use.}
\end{table}

\subsection{Activation Extraction}
\label{sec:extraction}

For each model--problem pair, we extract residual stream activations (the
output of each transformer block, post-MLP and post-residual connection) at
all layers and at a set of semantically meaningful token positions.  For
standard models, we extract at two positions: \textbf{P0} (last token of the
prompt, before generation begins) and \textbf{P2} (the answer token).
For reasoning models, we additionally extract at positions within the
thinking trace: \textbf{T0} (first token after \texttt{<think>}),
\textbf{T25/T50/T75} (25\%/50\%/75\% through the thinking trace), and
\textbf{Tend} (last token before \texttt{</think>}).

Activations are stored in BFloat16 in a shared HDF5 cache following a
fixed schema, along with per-head attention entropy metrics computed
incrementally during generation (we never materialize full attention
matrices).  The complete schema is specified in our data contract
documentation.  Total storage for the full benchmark across all four models
is approximately 4\,GB.

\subsection{Workstream 1: Linear Probes}
\label{sec:probes}

We train $\ell_2$-regularized logistic regression probes
(\texttt{LogisticRegressionCV} with 5-fold stratified cross-validation,
stratified by task category and target label) on residual stream activations
at each layer and token position.  We probe for four targets:

\begin{enumerate}[leftmargin=*,itemsep=1pt]
  \item \textbf{Task type}: conflict vs.\ no-conflict condition.
  \item \textbf{Correctness}: whether the model will answer correctly.
  \item \textbf{Bias susceptibility}: whether the model will produce the
  specific S1-lure answer (on conflict items only).
  \item \textbf{Processing mode}: within matched pairs, distinguishing the
  conflict from the no-conflict version of structurally identical problems.
\end{enumerate}

\paragraph{Controls.}  We apply three mandatory controls:
(i)~\emph{Hewitt--Liang control tasks} \citep{hewitt2019control}: probes
trained on randomly permuted labels establish a floor for probe
expressiveness.  Selectivity (true accuracy minus control accuracy) below
5 percentage points indicates the signal reflects probe capacity, not
representational structure.
(ii)~\emph{Permutation tests}: 1000 label shuffles per layer to establish
significance thresholds, using the North et al.\ correction
\citep{north2002note}.
(iii)~\emph{Leave-one-category-out}: training on 6 of 7 task categories and
testing on the held-out category to assess whether probes learn a
domain-general S1/S2 signal rather than category-specific features.

We report ROC-AUC as the primary metric with bootstrap 95\% confidence
intervals (1000 resamples).  We additionally train 2-layer MLP probes
(hidden dimension 256, ReLU activation) as a nonlinearity check and apply
Contrast-Consistent Search \citep[CCS;][]{burns2022discovering} as an
unsupervised corroboration for the correctness target.

\subsection{Workstream 2: SAE Feature Analysis}
\label{sec:sae}

We use pre-trained sparse autoencoders (Llama Scope for Llama-based models
\citep{he2024llama}; Gemma Scope for Gemma \citep{lieberum2024gemma}) to
decompose residual stream activations into sparse feature vectors.  For each
SAE feature, we compute the mean activation on conflict versus no-conflict
items and test for differential activation via Mann--Whitney~$U$ tests,
corrected for multiple comparisons using the Benjamini--Hochberg procedure
at $q = 0.05$ \citep{benjamini1995controlling}.

\paragraph{Visualization.}  We construct volcano plots (log fold-change vs.\
$-\log_{10} p$) to identify features with both large effect sizes and
statistical significance.  Features in the upper corners---large differential
activation between conditions---are candidate ``deliberation features'' and
``heuristic features.''

\paragraph{Falsification.}  Following \citet{ma2026falsification}, who
showed that 45--90\% of claimed SAE reasoning features are spurious
token-level artifacts, we apply their falsification protocol to every
candidate feature.  For each feature, we identify its top-3 activating tokens,
inject those tokens into 100 random non-cognitive-bias texts, and test whether
the feature still activates.  Features that survive this filter are classified
as genuine processing-mode features; we report falsification rates for all
candidates.

\subsection{Workstream 3: Attention Entropy}
\label{sec:attention}

We quantify attention focus via Shannon entropy
$H_{l,h,t} = -\sum_i a_i \log_2 a_i$ and the Gini coefficient of attention
weights, computed per-head per-position during generation.  We normalize
entropy by $\log_2 t$ (where $t$ is the current sequence length) to control
for the mechanical effect of longer sequences on entropy.

For each head, we test whether the attention metric distribution differs
between conflict and no-conflict items using Mann--Whitney~$U$ tests,
correcting across all heads $\times$ metrics via BH-FDR.  We classify heads
as \emph{S2-specialized} if they show significantly higher entropy on conflict
items in $\geq 3$ of 5 metrics with effect size $|r_{\text{rb}}| \geq 0.3$,
and \emph{S1-specialized} for the opposite pattern.

\paragraph{GQA handling.}  Because Llama uses grouped-query attention (32
query heads sharing 8 KV heads in groups of~4), heads within the same KV
group are not statistically independent.  We report results at both the
query-head and KV-group granularity, with the latter being the conservative
estimate.

\paragraph{Gemma-2 sliding window.}  Gemma-2's odd-numbered layers use a
4096-token sliding window for attention.  We analyze odd and even layers
separately and never pool them.

\subsection{Workstream 4: Representational Geometry}
\label{sec:geometry}

We analyze the geometric structure of activation point clouds corresponding to
S1-like and S2-like processing conditions.

\paragraph{Clustering.}  At each layer, we compute the \emph{cosine
silhouette score} between conflict and no-conflict activation clusters.  We
first reduce dimensionality to 50--100 principal components to avoid the
$d \gg N$ pitfall (with $d = 4096$ and $N \approx 284$, any two random
classes are linearly separable by Cover's theorem).  All silhouette scores
are compared against a permutation baseline (10,000 label shuffles).

\paragraph{Cross-model comparison.}  We compute layer-matched Centered Kernel
Alignment \citep[CKA;][]{kornblith2019similarity} between
Llama-3.1-8B-Instruct and R1-Distill-Llama-8B on conflict and no-conflict
items separately.  We predict that CKA diverges more between the two models
on conflict items (where reasoning distillation should produce the largest
representational changes) than on no-conflict items.

\paragraph{Intrinsic dimensionality.}  We estimate the intrinsic
dimensionality of S1-like and S2-like activation point clouds using the
Two-NN estimator \citep{facco2017estimating}.  Lower intrinsic
dimensionality for S2-like activations would suggest a more stereotyped
deliberation circuit; higher dimensionality would suggest diverse processing
strategies.

\paragraph{Random projection baseline.}  For any UMAP visualization, we
verify that apparent cluster structure survives projection to 2D via random
Gaussian matrices (100 draws).  If clusters appear in random projections,
the structure is genuine; otherwise, UMAP may be hallucinating separation.

\subsection{Workstream 5: Causal Interventions}
\label{sec:causal}

To move from correlation to causation, we intervene on SAE features and
measure behavioral effects.

\paragraph{Feature steering.}  For each candidate deliberation feature
identified in Workstream~2 (post-falsification), we perform activation
steering: we add $\alpha \cdot \hat{f}$ to the residual stream at the
feature's peak layer during inference, where $\hat{f}$ is the SAE decoder
direction and $\alpha$ is a steering coefficient.  We vary $\alpha$ over
$\{0.5, 1.0, 2.0, 3.0, 5.0\}$ to obtain dose-response curves.

\paragraph{Controls.}  We compare against (i)~steering along a randomly
sampled direction in activation space (same norm), which should produce
negligible behavioral change if the identified features are causally
specific; and (ii)~ablating the candidate features (setting their
activation to zero), which should shift behavior in the opposite direction
from amplification.

\paragraph{Metrics.}  We measure the change in (i)~probability of the
correct answer, (ii)~probability of the S1-lure answer, and (iii)~general
capability on a held-out set of non-bias tasks (to check for side effects).
All experiments use 3 random seeds.
